\chapter{Symbolic Views and Dialects}

The ISAR architecture supports multiple symbolic views or ``dialects'' over the representation-free substrate. Each view compiles to the substrate and decodes back, satisfying a preservation law. In this chapter, we formalize several dialects: the Lambda calculus fragment, ZFC hereditarily finite sets, term rewriting systems, bytecode compilation, and Barker's Iota combinator.

\section{The Lambda Calculus Fragment}

The bracket abstraction compiler and simulation properties in this section are mechanically verified in \url{src/ISAR/LambdaFragment.lean}. We define the standard $\lambda$-calculus with de Bruijn indices.

\begin{definition}[Lambda Terms]
\lean{ISAR.LTerm}
Lambda terms with variables (represented by natural number de Bruijn indices), applications, and abstractions:
\begin{align*}
t, u \in \text{LTerm} &::= \mathbf{var}\ n \mid \mathbf{app}\ t\ u \mid \mathbf{lam}\ t
\end{align*}
\end{definition}

\begin{definition}[Lambda Step]
\lean{ISAR.LStep}
Single-step $\beta$-reduction on \texttt{LTerm} using de Bruijn shifts and substitutions.
\end{definition}

\begin{definition}[Lambda Compiler]
\lean{ISAR.compile}
\uses{ISAR.LTerm, ISAR.ITerm}
Compiles lambda terms into the ISAR substrate via bracket abstraction:
$$\text{compile} : \text{LTerm} \to \text{ITerm}$$
\end{definition}

\begin{theorem}[Compiler Simulation]
\lean{ISAR.compile_simulates_step, ISAR.compile_simulates_red}
\uses{ISAR.compile, ISAR.LStep}
The compiler faithfully simulates lambda reductions in the ISAR kernel:
$$\forall t, u \in \text{LTerm}, \quad t \to_\beta u \implies \text{compile } t \to^* \text{compile } u$$
\end{theorem}

\section{Hereditarily Finite Sets and ZFC}

The set-theoretic encoding and relative ZFC interpretation verified in this section are mechanically verified in \url{src/ISAR/HFSet.lean}, \url{src/ISAR/HFSetEncoding.lean}, and \url{src/ISAR/HFSetSemantics.lean}. This relative interpretation is strictly limited to the Hereditarily Finite (HF) set fragment (ZFC without the axiom of infinity). We formalize HF Sets, their encoding, and show they satisfy set-theoretic axioms.

\begin{definition}[HF Sets]
\lean{HF}
The type of hereditarily finite sets, built inductively from the empty set:
$$x, y \in \text{HF} ::= \emptyset \mid x \cup \{y\}$$
\end{definition}

\begin{definition}[Extensional Equality]
\lean{ExtEq}
\uses{HF}
Set-theoretic extensional equivalence ($x =_{ext} y$) defined via Ackermann's bijective encoding `toNat`:
$$x =_{ext} y \iff \text{ExtEq } x\ y$$
\end{definition}

\begin{definition}[HF Set Encoding]
\lean{ISAR.HF_encode}
\uses{HF, ISAR.InvariantLayer}
Encodes HF sets into substrate invariant layers:
$$\text{HF\_encode} : \text{HF} \to \text{InvariantLayer}$$
\end{definition}

\begin{theorem}[HF Encoding Correctness]
\lean{ISAR.decode_layer_HF_encode, ISAR.HF_encode_decode_layer}
\uses{ISAR.HF_encode}
The encoding is sound and invertible:
$$\forall c \in \text{HF}, \quad \text{decode\_layer}(\text{HF\_encode } c) = c$$
\end{theorem}

We package HF Sets as a semantic kernel.

\begin{definition}[HF Set Kernel]
\lean{ISAR.HF_Kernel}
\uses{ISAR.Kernel, HF}
The set-theoretic admissible semantic kernel HF\_Kernel.
\end{definition}

\begin{theorem}[ZFC Interpretation Factorization]
\lean{ISAR.HF_Kernel_factorization}
\uses{ISAR.HF_Kernel, ISAR.ISAR_Kernel_terminal}
The HF set kernel factors uniquely through \texttt{ISAR\_Kernel}:
$$\forall f : \text{KernelHom HF\_Kernel ISAR\_Kernel}, \quad f(c) \sim_{op} \text{encode\_raw } c$$
This faithful factorization proves that the ZFC HF set fragment faithfully interprets into the ISAR substrate, satisfying empty set, pairing, and union axioms.

Syntactically distinct terms representing the same set (e.g., $\{a, b\}$ and $\{b, a\}$) project to the exact same element in the `InvariantLayer` quotient, demonstrating extensional equivalence.
\end{theorem}

\section{Term Rewriting, Bytecode, and Iota Views}

The SKI TRS view, stack bytecode compiler, and Barker's iota dialect are verified in \url{src/ISAR/TRSView.lean}, \url{src/ISAR/BytecodeView.lean}, and \url{src/ISAR/IotaView.lean} respectively. We define dialects for other standard computational models.

\begin{definition}[SKI TRS Dialect]
\lean{ISAR.TRS_Dialect}
\uses{ISAR.Dialect}
The term rewriting dialect modeling pure SKI combinator terms \texttt{TTerm}.
\end{definition}

\begin{definition}[Bytecode Dialect]
\lean{ISAR.Bytecode_Dialect}
\uses{ISAR.Dialect}
The dialect representing stack-based virtual machine instruction bytecode \texttt{Instruction} compiled and decompiled.
\end{definition}

\begin{definition}[Barker's Iota Dialect]
\lean{ISAR.Iota_Dialect}
\uses{ISAR.Dialect}
The dialect modeling the single-combinator \texttt{IotaTerm} utilizing Barker's universal combinator $\iota = \lambda x. x\ S\ K$.
\end{definition}

In Barker's iota combinator view, iterating iota-application produces the closed orbit:
$$ I \to A \to K \to S \to X \to I $$
where:
\begin{itemize}
    \item $I$ = identity ($I x = x$)
    \item $A$ = flip constant ($A x y = y$)
    \item $K$ = constant ($K x y = x$)
    \item $S$ = substitution ($S f g x = f x (g x)$)
    \item $X$ = self-application kernel
\end{itemize}
These four distinct active combinators map directly onto the four matrices ($I, R, A, S$) of the ISAR substrate, proving the ontological completeness of the 4-carrier representation.
